\section{Estimates and the Gantt Chart}
\label{sec:estimates-and-the-gantt-chart}

\begin{table}[H]
  \centering
  \begin{tabular}{@{}lcc@{}}
    \toprule
    Task & Estimated Hours & Output \\
    \midrule
    Analyze reference structure & 4 & Formatting notes \\
    Create LaTeX scaffold & 6 & Reusable project template \\
    Validate build chain & 2 & Working PDF generation \\
    \bottomrule
  \end{tabular}
  \caption{Illustrative task summary. \textit{Source: Compiled by the author.}}
  \label{tab:task-summary}
\end{table}

\begin{landscape}
\begin{figure}[p]
  \centering
  \begin{ganttchart}[
    x unit=0.7cm,
    y unit chart=0.8cm,
    vgrid,
    hgrid,
    title height=1,
    bar/.style={fill=upcblue!60},
    title/.style={fill=upcblue!10}
  ]{1}{6}
    \gantttitle{Illustrative Work Plan}{6} \\
    \gantttitlelist{1,...,6}{1} \\
    \ganttbar{Structure analysis}{1}{2} \\
    \ganttbar{Project scaffolding}{3}{5} \\
    \ganttbar{Compilation validation}{6}{6}
  \end{ganttchart}
  \caption{Illustrative Gantt chart included to preserve the same planning-oriented package stack as the reference report. \textit{Source: Compiled by the author.}}
  \label{fig:gantt-chart}
\end{figure}
\end{landscape}
